\chapter*{序論}
\addcontentsline{toc}{chapter}{序論}
\VerdierPhysicalPageAnchor{0018}
\VerdierSemanticAnchor{intro}

\VerdierSemanticAnchor{intro-opening}
本論文では、超ホモロジーの形式体系を提案する。

\VerdierSemanticAnchor{intro-1}
\noindent\textbf{1.}\quad
$\mathcal A$ と $\mathcal B$ を二つのアーベル圏、
$F:\mathcal A\to\mathcal B$ を加法的関手とする。$\mathcal A$ が十分な
入射対象をもつとき、$F$ の右導来関手の構成は今日では古典的である [1]：
$\mathcal A$ の各対象 $X$ は、「ホモトピーを除いて」一意な入射分解
$X\xrightarrow{\varepsilon}I^{\bullet}(X)$ をもつ。したがって、複体
$F(I^{\bullet}(X))$ のコホモロジー対象は $X$ によって標準同型を除いて
定まり、$X$ に関手的に依存する。こうして得られる関手が $F$ の右導来
関手である。

\VerdierSemanticAnchor{intro-1-homotopy-quasi}
このように、$\mathcal A$ の各対象 $X$ に、$\mathcal B$ の対象からなる
「ホモトピーを除いて」定まる複体を関手的に対応させ、そのコホモロジーを
とる。しかし通常のホモロジー代数の構成が与えるのは、ホモトピーを除いて
定まる複体ではなく、多くの場合\emph{擬同型を除いて}定まる複体である
（複体の射がコホモロジー対象上で同型を誘導するとき、その射を擬同型と
いう）。この点を説明するため、二つの例を挙げよう。

\VerdierSemanticAnchor{intro-1-f-acyclic}
まず、$F$ の導来関手を調べるには、入射分解より一般の分解
$X\xrightarrow{\varepsilon'}\mathcal R(X)$、すなわち関手 $F$ に関して
非輪状な対象による分解を考えるべきである（層の理論ではフラスク層または
軟層による分解 [23]、加群の理論では平坦加群による分解など）。これらの
分解は、もはやホモトピーを除いて一意とは限らない。しかし、$F$-非輪状
分解 $X\xrightarrow{\varepsilon'}\mathcal R(X)$ と入射分解
$X\xrightarrow{\varepsilon}I^{\bullet}(X)$ が与えられると、ホモトピーを
除いて一意な分解の射、すなわち次頁の図式を可換にする複体の射
$\rho:\mathcal R(X)\to I^{\bullet}(X)$ が存在する：
